\documentclass[11pt]{article}

\usepackage[margin=1in]{geometry}
\usepackage{amsmath}
\usepackage{amssymb}
\usepackage{amsthm}
\usepackage{booktabs}
\usepackage{enumitem}
\usepackage{float}
\usepackage[T1]{fontenc}
\usepackage{lmodern}
\usepackage{microtype}
\usepackage[numbers,sort&compress]{natbib}
\usepackage{xcolor}
\usepackage{hyperref}
\usepackage{tabularx}
\usepackage{url}

\hypersetup{
  colorlinks=true,
  linkcolor=blue!45!black,
  citecolor=blue!45!black,
  urlcolor=blue!45!black,
  pdftitle={VOLE-ARC: Scoped Anonymous Rate-Limited Credentials},
  pdfsubject={Protocol and implementation note}
}

\setlist[itemize,enumerate]{nosep,leftmargin=*}
\setcounter{tocdepth}{1}
\setlength{\parindent}{0pt}
\setlength{\parskip}{0.55em}
\raggedbottom

\newcommand{\concat}{\,\|\,}
\newcommand{\F}{\mathbb{F}_{16}}
\newcommand{\LE}[1]{\operatorname{LE}_{#1}}
\newcommand{\XOF}{\operatorname{SHAKE256}}
\newcommand{\Prefix}{\operatorname{Prefix}}
\newcommand{\pack}{\operatorname{pack16}}
\newcommand{\PoK}{\operatorname{PoK}}
\newcommand{\Store}{\ensuremath{\mathsf{Store}}}
\newcommand{\domain}[1]{\texttt{#1}}
\newcommand{\Adv}{\operatorname{Adv}}
\newcommand{\Bad}{\mathsf{Bad}}
\newcommand{\Issued}{\mathcal{I}}

\newtheorem{definition}{Definition}
\newtheorem{theorem}{Theorem}
\newtheorem{lemma}{Lemma}
\theoremstyle{remark}
\newtheorem{remark}{Remark}

\title{\textbf{VOLE-ARC: Scoped Anonymous Rate-Limited Credentials}}
\author{Protocol and implementation note}
\date{Implementation snapshot 0.1.0 -- 28 July 2026}

\begin{document}
\maketitle

\begin{abstract}
VOLE-ARC is an experimental anonymous credential with a separate presentation
allowance for each public context. A client obtains one hidden credential,
then proves possession while deriving a deterministic tag from a secret,
context, and hidden counter. A verifier accepts at most one copy of each tag.
The counter range therefore bounds accepted presentations per credential and
context. The construction follows the scoped rate-limit pattern of the
Privacy Pass ARC work, but replaces its P-256 credential and Schnorr proof
with signer-salted MAYO authentication and a VOLE-in-the-head proof of the
complete relation. This note specifies the implemented protocol and proves a
conditional scoped-accounting theorem. It also reduces the signer-salted
wrapper to an explicit auxiliary-preimage random-target assumption in the
classical ideal-XOF model. The proof does not establish that assumption for
MAYO, prove the pinned proof backend extractable or zero knowledge, or cover
fixed Keccak, quantum oracle access, side channels, or deployment failures.
VOLE-ARC is a research prototype, not a production scheme or an interoperable
Privacy Pass ciphersuite.
\end{abstract}

\tableofcontents
\newpage

\section{Scope and status}

The Privacy Pass ARC drafts let one credential produce a fixed number of
unlinkable tokens under an agreed presentation
context~\cite{arc-crypto,arc-protocol}. VOLE-ARC keeps that structure:

\begin{enumerate}
  \item issue a hidden, long-lived credential;
  \item derive a tag from a credential secret, a public presentation scope,
        and a hidden nonce;
  \item prove that the nonce is below a public limit; and
  \item reject a tag that was accepted before.
\end{enumerate}

The cryptography and encoding are different. ARC(P-256) uses a
keyed-verification credential, a group-element tag, and an integrated Schnorr
range proof. VOLE-ARC uses a MAYO preimage as a publicly verifiable
authenticator and proves the MAYO, SHAKE256, and range relations with
VOLE-in-the-head~\cite{mayo-round2,vole-in-the-head}. Its tag is 32 bytes.
Its local envelope starts with \domain{VARC}, not Privacy Pass token type
\domain{0xE5AC}. It is not wire-compatible with ARC(P-256), and it does not
specify an HTTP binding.

This document describes the Rust implementation at the repository's initial
snapshot. Section~\ref{sec:security} proves the protocol's accounting
composition under explicit assumptions and identifies the assumptions that
remain unproved for the implementation.

\section{Protocol specification}
\label{sec:protocol}

\subsection{Notation and roles}

Byte strings are concatenated with $\concat$. $\LE{w}(x)$ encodes integer
$x$ as $w$ little-endian bytes. All lengths in transcript framing are
unsigned 64-bit little-endian integers. $\XOF(x,n)$ denotes the first $n$
bytes of SHAKE256 on $x$~\cite{fips-202}. For a domain string $d$, define
\[
  H_d(x_1,\ldots,x_r) =
  \XOF\!\left(d \concat
    \LE{8}(|x_1|) \concat x_1 \concat \cdots \concat
    \LE{8}(|x_r|) \concat x_r, 32\right).
\]
Domain strings include the ASCII suffix \domain{/v1}. This suffix is part of
the transcript, not prose versioning.

The three roles are:

\begin{itemize}
  \item \textbf{Client.} Creates an issuance proof, retains the credential
        secret and
        counter map, and constructs presentations.
  \item \textbf{Issuer.} Verifies issuance proofs and uses a MAYO trapdoor to
        authenticate a client commitment.
  \item \textbf{Verifier.} Reconstructs policy-selected public inputs,
        verifies a
        presentation proof, and atomically consumes its tag.
\end{itemize}

The verifier can be a relying party distinct from the issuer because
presentation verification is public. The issuer public key must still be
authenticated out of band.

\subsection{Suites and keys}

Let $P^*:\F^{kn}\rightarrow\F^m$ be the whipped MAYO public map. The issuer
holds its trapdoor and publishes the expanded map, an application/key-epoch
label $A$, and one VOLE profile. The profile selects the proof tree geometry
through two parameters, denoted $(\tau,\kappa_V)$. The implementation
supports only the two level-1 MAYO parameter sets shown below.

\begin{table}[H]
  \centering
  \caption{Supported combined suites. MAYO values follow the Round 2
  specification~\cite{mayo-round2}.}
  \begin{tabular}{lrrrr}
    \toprule
    Suite & $n$ & $m$ & $o$ & $k$ \\
    \midrule
    MAYO1 & 86 & 78 & 8 & 10 \\
    MAYO2 (default) & 81 & 64 & 17 & 4 \\
    \bottomrule
  \end{tabular}
\end{table}

The implementation first hashes the expanded public quadratic forms in their
canonical visitation order:
\begin{align*}
  h_{\mathrm{pk}} = \XOF(&
    \domain{VOLE-ARC/mayo-public-key/v1} \concat
    \LE{8}(n) \concat \LE{8}(m) \concat \LE{8}(o) \concat \LE{8}(k) \\
    &{}\concat \operatorname{forms}(P^*), 32).
\end{align*}
It then derives the issuer context
\begin{align*}
 I = \XOF(&\domain{VOLE-ARC/issuer-context/v1} \concat
   \LE{8}(|A|) \concat A \concat h_{\mathrm{pk}} \concat
   \operatorname{name}(P) \\
   &{}\concat \LE{8}(32) \concat \LE{8}(\tau) \concat
   \LE{8}(\kappa_V),32).
\end{align*}
The value 32 fixes the presentation-counter width. The domain version, suite
name, and profile parameters keep keys for distinct relations in distinct
protocol contexts.

\subsection{Application contexts}

Four public values separate credential policy, rate-limit policy, and request
binding:

\begin{table}[H]
  \centering
  \caption{Public application inputs.}
  \begin{tabularx}{\textwidth}{@{}lXl@{}}
    \toprule
    Value & Meaning & Expected lifetime \\
    \midrule
    $CC$ & Credential context, such as an attestation class or credential
      epoch & Credential lifetime \\
    $PC$ & Presentation context naming a relying party, purpose, and optional
      epoch & One rate-limit bucket \\
    $B$ & Digest of a challenge or protected request & One request \\
    $L$ & Nonzero, public presentation limit & Server policy \\
    \bottomrule
  \end{tabularx}
\end{table}

For arbitrary application bytes $a$ and request bytes $b$:
\begin{align*}
  CC &= H_{\domain{VOLE-ARC/credential-context/v1}}(a),\\
  B  &= H_{\domain{VOLE-ARC/presentation-binding/v1}}(b).
\end{align*}
A generic bucket is
\[
  PC = H_{\domain{VOLE-ARC/presentation-context/v1}}(p).
\]
The convenience constructor for relying party $R$, purpose $U$, and optional
epoch $E$ instead computes
\[
  PC = H_{\domain{VOLE-ARC/scoped-presentation/v1}}(R,U,e,E),
\]
where $e$ is the one-byte value 1 when an epoch is present and 0 otherwise;
$E$ is empty when absent. Framing distinguishes every component, including
an absent epoch from a present empty epoch.

Two fixed-width effective scopes keep the in-circuit SHAKE messages within
one 136-byte rate block:
\begin{align*}
  CS &= \XOF(\domain{VOLE-ARC/credential-scope/v1}
             \concat I \concat CC,32),\\
  PS &= \XOF(\domain{VOLE-ARC/presentation-scope/v1}
             \concat I \concat CC \concat PC,32).
\end{align*}
Neither $B$ nor $L$ appears in $PS$. A new request does not reset a bucket.
Raising $L$ extends the existing nonce range; it does not create a second
allowance. The verifier, not the client, must select $PC$ from trusted policy.

\subsection{Credential issuance}

The client samples independent uniform values
\[
  k \leftarrow {\{0,1\}}^{256},
  \qquad
  \rho \leftarrow {\{0,1\}}^{256}.
\]
Here $k$ is the scoped-tag key and $\rho$ hides its credential commitment.
It computes
\[
  C = \Prefix_m\!\left(
      \XOF(\domain{VOLE-ARC/credential/v1}
           \concat CS \concat k \concat \rho,\lceil m/2\rceil)
      \right) \in \F^m .
\]
$\Prefix_m$ reads the low nibble and then the high nibble of each byte,
discarding the final high nibble if $m$ is odd.

The client sends $(C,\pi_{\mathrm{issue}})$, where the VOLE-in-the-head proof
establishes
\[
  \pi_{\mathrm{issue}} =
  \PoK\{(k,\rho): C =
  \operatorname{CredentialHash}(CS,k,\rho)\}.
\]
The Fiat-Shamir statement is the byte string
\[
  \domain{VOLE-ARC/issue-statement/v1} \concat
  \LE{8}(|\operatorname{name}(P)|) \concat \operatorname{name}(P) \concat
  I \concat CC \concat \LE{8}(m) \concat C_{\mathrm{bytes}}.
\]
Each field element in $C_{\mathrm{bytes}}$ occupies one byte in the statement.

The issuer verifies the proof under its policy-selected $CC$. Only after
acceptance does it sample a fresh 32-byte salt $\zeta$. It computes
\begin{align*}
  T = \Prefix_m\!\left(
      \XOF(\domain{VOLE-ARC/signed/v1}
           \concat \pack(C) \concat \zeta,\lceil m/2\rceil)
      \right)
\end{align*}
and samples a MAYO preimage
\[
  \sigma \leftarrow \operatorname{SPre}(sk,T)
  \quad\text{such that}\quad P^*(\sigma)=T.
\]
The response is $(\sigma,\zeta)$. The client checks $P^*(\sigma)=T$ against
its retained $C$ before storing
\[
  \mathsf{cred}=(I,CC,\sigma,k,\rho,\zeta,\mathsf{next}=\varnothing).
\]

The signer chooses $\zeta$ after proof verification so the client cannot
choose the exact MAYO target. Reissuing a response for the same request
changes $(\sigma,\zeta)$ but not $(k,\rho)$; all such responses remain one
tag lineage. This ordering is a design condition, not by itself a
one-more-preimage proof.

\subsection{Presentation}

A challenge is $(CC,PC,L,B)$ with $1\leq L\leq 2^{32}-1$. The client requires
that the challenge's $CC$ match the credential. It reads
$i=\mathsf{next}[PC]$, using zero when the bucket is new, and rejects when
$i\geq L$. It derives
\[
  t = \XOF(\domain{VOLE-ARC/tag/v1}
           \concat k \concat PS \concat \LE{4}(i),32).
\]

The presentation proof establishes knowledge of
$(\sigma,k,\rho,\zeta,i)$ satisfying
\begin{align}
 C &= \operatorname{CredentialHash}(CS,k,\rho), \label{eq:cred-hash}\\
 T &= \operatorname{SignedCredentialHash}(C,\zeta), \label{eq:signed-hash}\\
 P^*(\sigma) &= T, \label{eq:mayo}\\
 t &= \operatorname{TagHash}(k,PS,i), \label{eq:tag}\\
 0 &\leq i < L. \label{eq:range}
\end{align}
The commitment $C$ and signed target $T$ are hidden intermediate wires.
The Fiat-Shamir statement contains
\[
\begin{split}
 &\domain{VOLE-ARC/presentation-statement/v1} \concat
 \LE{8}(|\operatorname{name}(P)|) \concat \operatorname{name}(P) \concat I\\
 &\qquad{}\concat CC \concat PC \concat \LE{4}(L) \concat B \concat t.
\end{split}
\]

The circuit bit-decomposes $i$ in little-endian order. It constrains the
carry chain for the 32-bit addition $i+\neg L+1$. With initial carry
$c_{-1}=1$, each $c_j$ is the majority bit of
$(i_j,{\left(\neg L\right)}_j,c_{j-1})$. The final carry is one exactly when
$i\geq L$, so the circuit constrains $c_{31}=0$.

After proof construction succeeds, the client stores
$\mathsf{next}[PC]\leftarrow i+1$ and returns $(t,\pi_{\mathrm{present}})$.
The nonce and credential are not transmitted.

\subsection{Verification and scoped accounting}

The verifier reconstructs $CS$, $PS$, the circuit, and the Fiat-Shamir
statement from its authenticated public key and challenge policy. It verifies
$\pi_{\mathrm{present}}$. On success it performs one operation:
\[
  \Store.\operatorname{insert\_if\_absent}(PS,t).
\]
The presentation is accepted only when this operation inserts the first
durable copy.

Theorem~\ref{thm:accounting} gives the exact accounting argument. In brief,
extraction assigns each accepted presentation to an authenticated credential
lineage and an integer $i<L$. The tag is deterministic for that lineage,
$PS$, and $i$, so the store can accept each such pair at most once. With $q$
issued lineages and a fixed limit $L$, at most $qL$ presentations are
accepted. If limits change without changing $PS$, the bound is
$qL_{\max}$, where $L_{\max}$ is the largest accepted limit.

Tag collision resistance is not needed for this upper bound. A collision
between distinct tag inputs makes the store merge uses, reducing capacity or
causing denial of service. The load-bearing binding property is instead that
one authenticated commitment cannot open to two credential inputs and hence
two tag keys.

\section{State, encoding, and failure boundaries}

\subsection{Client state}

The counter map is privacy-critical and is not authenticated by
Equations~\eqref{eq:cred-hash}--\eqref{eq:mayo}. The canonical credential
encoding can detect a malformed or cryptographically invalid credential, but
it cannot detect a removed counter, a lower counter, or an older authentic
snapshot. Storage therefore needs confidentiality, integrity, and monotonic
rollback protection.

The required send order is:
\begin{enumerate}
  \item construct the presentation;
  \item atomically persist the incremented credential state; and
  \item transmit the presentation.
\end{enumerate}
A crash after step 2 can consume capacity. A rollback can repeat a public tag,
which links the retry and should be rejected. It does not add capacity while
the verifier store remains intact. Concurrent copies of one credential can
also emit the same tag.

The implementation retains at most 4096 presentation-context counters. It
fails when a new context would exceed that bound. Silent eviction would reset
a counter and is not permitted.

\subsection{Verifier state}

\Store must be linearizable across verifier replicas and durable before it
reports a successful insertion. Its recovery point must be at least as recent
as the protected application state. The protected action and tag consumption
need transaction semantics appropriate to the application. The included
in-memory store provides only single-process test behavior.

Proof verification occurs before insertion. Invalid clients can therefore
force proof-verification work. Wire-size limits do not replace admission,
concurrency, and CPU controls.

\subsection{Canonical envelope}

Every encoded artifact begins with this nine-byte header:
\[
  \domain{VARC} \concat
  \domain{01} \concat
  \operatorname{artifact} \concat
  \operatorname{suite} \concat
  \domain{0000}.
\]
The last two bytes are reserved and must be zero. Variable-length byte strings
use a four-byte little-endian length. Nibbles are packed low first. Decoders
reject a wrong artifact or suite, an unsupported version, nonzero reserved
bytes, noncanonical nibble padding, unsorted or duplicate counters, zero
stored counters, trailing bytes, and artifacts over 32 MiB.

\begin{table}[H]
  \centering
  \caption{Artifact identifiers and exposure.}
  \begin{tabularx}{\textwidth}{@{}clX@{}}
    \toprule
    ID & Artifact & Boundary \\
    \midrule
    1 & Public key & Public configuration and expanded MAYO map \\
    2 & Issuer key & Secret local state; unsealed \\
    3 & Issue request & Network artifact: packed $C$ and proof \\
    4 & Issue response & Network artifact: packed $\sigma$ and $\zeta$ \\
    5 & Pending issue & Secret local state; unsealed \\
    6 & Credential & Secret local state and counter map; unsealed \\
    7 & Presentation & Network artifact: tag and proof \\
    \bottomrule
  \end{tabularx}
\end{table}

The envelope is a local research format. It is not a Privacy Pass TLS
structure, and no independent implementation provides a byte-level
interoperability oracle.

\section{Parameters and measured cost}

The combined suite is limited by the VOLE layer, not only by MAYO. Both
profiles use lambda-128 VOLE parameters and 128-bit hidden proof-tree seeds.
Adding MAYO3 or MAYO5 without widening and reanalyzing the proof system would
misstate the security of the result. NIST selected MAYO to continue to the
third evaluation round in May 2026, but it remains a candidate rather than a
standard~\cite{nist-ir-8610}.

For MAYO2, $m=64$ makes $C$ and $T$ 256 bits. A generic collision has an
ideal cost near $2^{128}$ classically and $2^{256/3}\approx 2^{85.3}$
quantum queries under a BHT-style collision model. MAYO1 uses $m=78$, giving
312-bit values and corresponding ceilings near $2^{156}$ and $2^{104}$.
The degree-16 polynomial assertion contributes the bounded term
\[
  \frac{16+1}{2^{128}} = \frac{17}{2^{128}} \approx 2^{-123.91}.
\]
These are component ceilings, not a complete advantage bound. The proof
system's 128-bit hidden seeds also have a generic $2^{64}$ Grover query
ceiling.

Table~\ref{tab:bench} records one machine-specific balanced-profile snapshot.
It is useful for implementation planning, not cross-platform prediction.
The source benchmark uses Criterion with seeded fixtures and excludes network
and durable-database latency.

\begin{table}[H]
  \centering
  \caption{Balanced profile mean latency on an Apple M4 Pro, 28 July 2026.}
  \label{tab:bench}
  \begin{tabular}{lrr}
    \toprule
    Operation & MAYO2 & MAYO1 \\
    \midrule
    Issuance, end to end & 13.28 ms & 18.50 ms \\
    Presentation proof & 26.48 ms & 38.44 ms \\
    Presentation verification & 7.37 ms & 15.62 ms \\
    Presentation, end to end & 33.86 ms & 54.61 ms \\
    \bottomrule
  \end{tabular}
\end{table}

\begin{minipage}{\textwidth}
The balanced MAYO2 issue request is 55,286 bytes and its presentation is
138,550 bytes. The expanded public-key envelope is 106,329 bytes. These sizes
come largely from the generic VOLE proof, which checks three in-circuit
SHAKE256 evaluations, the whipped MAYO map, and the 32-bit comparison. They
explain why this prototype is slower and larger than a specialized
group-based ARC construction.
\end{minipage}

\section{Security status and open proof obligations}
\label{sec:security}

\subsection{Model}

Let $\mathcal{Y}=\F^m$ and $\mu=\log_2|\mathcal{Y}|=4m$. Write the three
secret-input functions as
\begin{align*}
 H_C(CS,k,\rho) &\in \mathcal{Y},\\
 H_S(C,\zeta) &\in \mathcal{Y},\\
 H_T(k,PS,i) &\in {\{0,1\}}^{256}.
\end{align*}
These abbreviations include the domains and encodings in
Section~\ref{sec:protocol}. The accounting theorem treats them as arbitrary
deterministic functions. It does not model SHAKE256 as random.

The proof system is used as an adaptive, multi-theorem non-interactive
argument of knowledge. Its extractor must return $(k,\rho)$ for every
accepted issuance proof before the issuer authenticates that commitment, and
$(\sigma,k,\rho,\zeta,i)$ for every accepted presentation before the store
authorizes the protected action. Extraction is straight-line and online:
rewinding an already stateful issuer or verifier is not part of the
experiment. Let $\epsilon_{\mathrm{ext}}$ bound the probability that any
accepted proof in the experiment lacks such an extracted witness satisfying
the exact relation and statement.

Each successful issuance response is entered in an extraction ledger
\[
  \Issued=\{(CS,C,k,\rho):
      C=H_C(CS,k,\rho)
      \text{ and the issuer returned an authenticator for }C\}.
\]
Repeated requests with the same extracted input make one lineage. The
following events cover the other ways an accepted presentation could escape
that ledger:
\begin{itemize}
  \item $\Bad_{\mathrm{auth}}$: an extracted tuple contains
        $P^*(\sigma)=H_S(C,\zeta)$ for a commitment $C$ that received no
        successful issuance response;
  \item $\Bad_{\mathrm{bind}}$: two distinct relevant inputs
        $(CS,k,\rho)\ne(CS',k',\rho')$ have the same $H_C$ output; and
  \item $\Bad_{\mathrm{store}}$: the store accepts the same $(PS,t)$ twice,
        loses a successful insertion before the protected action is
        irreversible, authorizes more than one protected action per
        insertion, or lets an action occur without a successful insertion.
\end{itemize}
Write their probabilities as $\epsilon_{\mathrm{auth}}$,
$\epsilon_{\mathrm{bind}}$, and $\epsilon_{\mathrm{store}}$. The store term is
an operational assumption, not a cryptographic advantage.

\begin{lemma}[Range relation]
\label{lem:range}
For $i,L\in\{0,\ldots,2^{32}-1\}$, with $L>0$, the circuit's final carry is
zero exactly when $i<L$.
\end{lemma}

\begin{proof}
The circuit adds $i+\neg L+1$ as 32-bit unsigned integers. Since
$\neg L=2^{32}-1-L$, the untruncated sum is $i+2^{32}-L$. It has a carry out
of bit 31 exactly when it is at least $2^{32}$, equivalently when $i\geq L$.
The circuit constrains that carry to zero.
\end{proof}

\subsection{Scoped-accounting theorem}

\begin{theorem}[Scoped accounting]
\label{thm:accounting}
Fix one authenticated issuer key and effective values $CS$ and $PS$. Suppose
at most $q$ distinct entries of $\Issued$ have credential scope $CS$. Let an
adaptive adversary submit presentations under challenges whose reconstructed
credential and presentation scopes are $CS$ and $PS$, whose bindings may
differ, and whose nonzero limits satisfy $L_j\leq L_{\max}$. If
$\mathsf{Acc}_{PS}$ is the number of protected actions authorized by first
insertions in that bucket, then
\[
 \Pr[\mathsf{Acc}_{PS}>qL_{\max}]
 \leq
 \epsilon_{\mathrm{ext}}+
 \epsilon_{\mathrm{auth}}+
 \epsilon_{\mathrm{bind}}+
 \epsilon_{\mathrm{store}}.
\]
\end{theorem}

\begin{proof}
Condition on the complement of all four bad events. Extract an accepted
presentation witness and set
\[
 C=H_C(CS,k,\rho),\qquad
 t=H_T(k,PS,i).
\]
Lemma~\ref{lem:range} gives $i<L_j\leq L_{\max}$. Since authentication does
not fail, $C$ is the commitment of an entry in $\Issued$. Since binding does
not fail, the extracted $(CS,k,\rho)$ equals that issued entry's input. Thus
every accepted presentation is assigned to a pair
\[
   (\text{issued lineage},i)
   \in \{1,\ldots,q\}\times\{0,\ldots,L_{\max}-1\}.
\]
Two accepted presentations assigned to the same pair have identical $k$,
$PS$, and $i$, hence the same deterministic tag $t$. The store cannot accept
both. The assignment is therefore injective, and its codomain has
$qL_{\max}$ elements. Taking a union bound over the excluded events proves
the result.
\end{proof}

\begin{remark}[What the theorem does not require]
The proof does not assume that $H_T$ is collision resistant. Distinct pairs
that collide under $H_T$ compete for one store entry, so a collision can only
reduce the number accepted. Tag pseudorandomness is required for privacy, not
for the upper bound. Honest client counter monotonicity is also unnecessary
for accounting: rollback repeats a pair and the store rejects it. Client
monotonicity instead prevents linkable retries and lost capacity.
\end{remark}

For a fixed limit $L$, the theorem gives the advertised $qL$ bound. When a
server first uses $L=5$ and later $L=10$ under the same $PS$, it gives ten
total positions per lineage, not five plus ten. The theorem counts
credentials, not people: issuance policy must enforce any intended
per-principal value of $q$.

The count above is by effective $CS$. To identify it with a named
$(I,CC)$ policy, the scope hash must not collide on two policy inputs.
Likewise, separate relying-party or epoch buckets require distinct effective
$PS$ values. A $CS$ collision can import a lineage from another named
credential policy, while a $PS$ collision merges store buckets. Neither
violates the theorem when $q$ and the bucket are defined by effective values,
but both invalidate the corresponding named-context claim.

\begin{remark}[Request binding]
The value $B$ is bound to the Fiat-Shamir statement but does not enter the
witness equations or $PS$. This is deliberate: changing a request must not
reset the bucket. The accounting theorem remains true even if a proof could
transfer between two $B$ values, because its tag still enters the same store.
The stronger claim that a captured proof verifies only for its original $B$
requires statement binding of the exact Fiat-Shamir transcript. Knowledge
soundness alone does not imply non-transferability between two true
statements.
\end{remark}

\subsection{Classical ideal-XOF bounds}

This subsection makes the hash terms and signer-salted authentication
assumption explicit. Each domain-separated SHAKE256 use is modeled as an
independent classical random oracle with the output length used by the
protocol. Key generation, credential secrets, nonces, salts, sampler coins,
and proof coins are assumed to use independent cryptographically secure
randomness.

If $Q_C$ distinct credential-hash inputs are evaluated across honest
generation, adversarial queries, extracted witnesses, and verification, then
\begin{equation}
\label{eq:binding-bound}
  \epsilon_{\mathrm{bind}}
  \leq \frac{\binom{Q_C}{2}}{2^\mu}.
\end{equation}
This count must include offline grinding; counting only the two openings
eventually presented would be invalid. For either 256-bit scope hash,
$Q_{\mathrm{scope}}$ distinct inputs collide with probability at most
$\binom{Q_{\mathrm{scope}}}{2}/2^{256}$.

The raw whipped MAYO map is not used through the standard MAYO signature
interface. The required property must therefore name the exact trapdoor
sampler.

\begin{definition}[Multi-target auxiliary-preimage one-wayness]
\label{def:mt-apow}
The $\mathsf{MT\text{-}APOW}_{P^*,\mathrm{SPre}}(q_A,q_T)$ game generates a
MAYO trapdoor and publishes $P^*$. An auxiliary-sample query chooses uniform
$Y\leftarrow\mathcal{Y}$ and runs the exact
$\mathrm{SPre}(sk,Y)$ algorithm, including its randomness, rejection
behavior, and possible failure. It returns $(Y,\sigma)$ on success and
$(Y,\bot)$ on failure. A target query returns an independent uniform
$Y\leftarrow\mathcal{Y}$. The adversary wins if it gives a preimage under
$P^*$ of either a target-query value or an auxiliary-sample value for which
it received $\bot$. Values may collide across queries; such collisions are
part of the game's advantage. Let
\[
 \Adv^{\mathsf{mt\text{-}apow}}_{P^*,\mathrm{SPre}}(q_A,q_T)
\]
be the maximum winning probability over efficient adversaries making at most
$q_A$ auxiliary-sample and $q_T$ target queries.
\end{definition}

This is a multi-target random one-wayness assumption in the presence of exact
signer samples. It is stronger and more specific than merely saying that
MAYO signatures are existentially unforgeable.

\begin{lemma}[Signer-salted wrapper]
\label{lem:salted}
Model $H_S$ as a classical random oracle. Give an adversary at most $q_I$
authentication attempts. On input $C$, the oracle samples
$\zeta\leftarrow{\{0,1\}}^{256}$ after receiving $C$, evaluates
$Y=H_S(C,\zeta)$, runs the exact $\mathrm{SPre}(sk,Y)$ algorithm, and records
$C$ only when it returns a response. If the adversary makes at most $Q_H$
other distinct $H_S$ queries, then its probability of outputting
$(C^*,\zeta^*,\sigma^*)$ such that
\[
 C^* \text{ was not recorded},\qquad
 P^*(\sigma^*)=H_S(C^*,\zeta^*)
\]
is at most
\begin{equation}
\label{eq:salted-bound}
 \Adv^{\mathsf{mt\text{-}apow}}_{P^*,\mathrm{SPre}}
      (q_I,Q_H+1)
 + \frac{q_I Q_H+\binom{q_I}{2}}{2^{256}}.
\end{equation}
\end{lemma}

\begin{proof}
Allowing authentication attempts on arbitrary $C$ only strengthens the real
issuer, which first requires a valid issuance proof.
Simulate $H_S$ lazily. A new adversarial hash query consumes one target query
from Definition~\ref{def:mt-apow}. On an authentication attempt, sample
$\zeta$. If $(C,\zeta)$ is fresh and distinct from earlier authentication
inputs, consume one auxiliary-sample query and program its uniform $Y$ at
that point. The returned sampler outcome has exactly the real distribution.

The simulation aborts if the adversary queried $(C,\zeta)$ before the
authentication attempt, or if an earlier attempt used the same pair. The
first event has probability at most $q_I Q_H/2^{256}$; the second at most
$\binom{q_I}{2}/2^{256}$. These are union bounds over fresh uniform salts.

If the final $(C^*,\zeta^*)$ was queried, a valid $\sigma^*$ inverts its
target-query value. If it was not queried, make one final target query and
test $\sigma^*$; revealing that target only after $\sigma^*$ is fixed cannot
weaken the reduction. A final label matching a successful authentication
label would have recorded the same $C^*$ and is excluded. A label from a
failed authentication is a winning failed auxiliary-sample value. Target
collisions with successful signer samples already count as wins in
Definition~\ref{def:mt-apow}, so no separate collision term is hidden.
\end{proof}

Combining Theorem~\ref{thm:accounting},
Equation~\eqref{eq:binding-bound}, and Lemma~\ref{lem:salted} gives the
classical ideal-XOF accounting bound
\begin{align}
\Pr[\mathsf{Acc}_{PS}>qL_{\max}]
\leq{}&
 \epsilon_{\mathrm{ext}}+\epsilon_{\mathrm{store}}
 + \frac{\binom{Q_C}{2}}{2^\mu} \notag\\
&+
 \Adv^{\mathsf{mt\text{-}apow}}_{P^*,\mathrm{SPre}}(q_I,Q_H+1)
 + \frac{q_I Q_H+\binom{q_I}{2}}{2^{256}}.
\label{eq:combined-bound}
\end{align}
Add the applicable 256-bit scope-collision term when $q$ is counted by named
application contexts rather than effective $CS$. In particular,
$\epsilon_{\mathrm{auth}}$ in Theorem~\ref{thm:accounting} is bounded by
Lemma~\ref{lem:salted}.

\subsection{Blind issuance and presentation privacy}

Accounting reveals equality exactly where needed: the same
$(k,PS,i)$ always reveals the same tag. Privacy can only be claimed for
histories with the same public metadata and the same equality pattern of
those hidden inputs. In particular, the issuer key, suite, profile, $CC$,
$PC$, $L$, $B$, message count, and network-visible behavior must be matched.

\begin{theorem}[Conditional transcript privacy]
\label{thm:privacy}
Consider $h$ honest, unexposed credentials whose $(k,\rho)$ pairs are sampled
independently and uniformly. Model $H_C$ and $H_T$ as independent classical
random oracles. Assume the issue and presentation arguments are adaptive
multi-theorem zero knowledge for the exact, possibly adversarially selected
public statements, with distinguishing advantage
$\epsilon_{\mathrm{zk}}$. A view containing their issuance transcripts and
presentations is indistinguishable from one in which:
\begin{enumerate}
  \item issue and presentation proofs are simulated;
  \item each honest credential commitment is a programmed uniform element of
        $\mathcal{Y}$; and
  \item each new honest tag input has a programmed uniform 256-bit output,
        while a repeated input repeats its output.
\end{enumerate}
All other public fields and adversarial messages are unchanged.
If the adversary makes $Q_C^{\mathrm{priv}}$ credential-hash queries and
$Q_T^{\mathrm{priv}}$ tag-hash queries, the distinguishing advantage is at
most
\begin{equation}
\label{eq:privacy-bound}
 \epsilon_{\mathrm{zk}}
 + \frac{hQ_C^{\mathrm{priv}}}{2^{512}}
 + \frac{hQ_T^{\mathrm{priv}}}{2^{256}}.
\end{equation}
\end{theorem}

\begin{proof}
Replace the proofs one at a time by the adaptive zero-knowledge simulator.
For each honest credential, sample its visible commitment uniformly and
program $H_C$ at the hidden $(CS,k,\rho)$ input. Until that exact 512-bit
secret pair is queried, the real and programmed views are identical. A union
bound over credentials and oracle queries gives the second term of
Equation~\eqref{eq:privacy-bound}.

Next sample each tag uniformly and program $H_T$ at
$(k,PS,i)$, reusing an existing value for a repeated input. Conditioned on
the earlier hybrids, each unexposed $k$ remains uniform. Even if $PS$ and
$i$ are known, a tag-oracle query hits one of the $h$ hidden keys with
probability at most $h/2^{256}$. The final union bound gives the third term.
\end{proof}

The ideal view makes the issuer's issuance-side transcript independent of the
tag key except for exhaustive hidden-input queries, and makes tags on distinct
inputs independent uniform strings. It therefore yields blind issuance and
issuance-to-presentation and pairwise-presentation unlinkability under the
theorem's matched-history conditions. A repeated tag, unusual public context,
different limit, rollback pattern, or network identifier remains linkable by
design or by the stated boundary. Accidental collisions between distinct
uniform tags have probability at most $\binom{n_T}{2}/2^{256}$ over $n_T$
honest tag inputs; they affect availability, not the accounting bound.

\subsection{Instantiation boundary}

The structural Theorem~\ref{thm:accounting} is independent of the random
oracle model, but its assumptions are not yet theorems about the pinned code.
The implementation pins MAYO and VOLE dependencies to VOLE-ACT revision
\domain{e3734dd}~\cite{vole-act}. A source pin fixes code; it does not prove:
\begin{itemize}
  \item adaptive multi-theorem straight-line extraction or zero knowledge for
        the exact degree-16 VOLE-in-the-head relations and Fiat-Shamir
        transcript;
  \item statement binding of that transcript to every public byte, including
        the request binding $B$;
  \item $\mathsf{MT\text{-}APOW}$ for the expanded MAYO1 or MAYO2 map with
        the exact bounded-retry sampler;
  \item a bridge from the oracle-augmented proof above to the circuit that
        evaluates fixed Keccak internally; or
  \item the same reductions for quantum superposition queries.
\end{itemize}
The degree-16 assertion term $17/2^{128}$ is one component of
$\epsilon_{\mathrm{ext}}$, not a complete value for it. Multi-user proof
attempts, vector-commitment security, VOLE consistency, Fiat-Shamir,
composition losses, and all oracle queries must also be counted.

Thus Equation~\eqref{eq:combined-bound} closes the protocol-level classical
composition only conditionally. It is not a standard-model proof for SHAKE256,
a QROM proof, a reduction of $\mathsf{MT\text{-}APOW}$ to the published MAYO
assumptions, or an audit of the implementation.

\subsection{Side channels and secret lifetime}

Locally owned secret buffers, circuit witnesses, Keccak states, and decoder
staging buffers are zeroized on a best-effort basis. Credential
authentication compares field encodings in constant time. Presentation
construction avoids a redundant native MAYO evaluation.

Those measures do not make the full implementation constant time. Dependency
allocations and copies are outside the wrapper's wiping guarantees. The
pinned prover contains secret-dependent branches in constraint bookkeeping
and VOLE mask construction. The counter map is a non-oblivious tree whose
allocation and lookup patterns can expose local usage history.

\begin{minipage}{\textwidth}
A fixed-vs-fixed timing diagnostic measured MAYO1 signing, MAYO2 signing, and
MAYO2 presentation proving. In one 2026-07-28 run, all reported Welch
statistics satisfied $|t|\leq 1.076$. This is a regression signal for one
compiler, machine, and input partition. It is not constant-time
certification. No machine-code audit, ctgrind run, power or EM analysis, fault
campaign, or traffic analysis has been completed.
\end{minipage}

\subsection{Deployment boundary}

Even a successful cryptographic proof would not supply the following:
\begin{itemize}
  \item policy that prevents one principal from obtaining extra credentials;
  \item canonical, server-controlled relying-party and epoch construction;
  \item authenticated public-key distribution, rotation, and revocation;
  \item rollback-resistant client storage and a replicated durable tag store;
  \item transaction coupling between tag insertion and the protected action;
  \item resistance to request flooding, network linkage, cookies, or logs; or
  \item compiler, hardware, RNG, and fault resistance.
\end{itemize}

Production use requires an independently reviewed reduction, independent
cryptanalysis, a second implementation with stable test vectors, sustained
fuzzing, side-channel and fault review, and crash-tested replicated state.
Unit tests, limited fuzz campaigns, source-level zeroization, and the timing
diagnostic are evidence about specific implementation paths only.

\section{Implementation conformance checks}

The repository makes review obligations executable where possible. One
canonical script drives Rust formatting, Clippy, unit tests with parallel and
sequential backends, warning-free Rustdoc, shell linting, paper linting and
build, fuzz-target build and lint, license/source policy, and advisory scans.
Tracked pre-commit and pre-push hooks invoke the quick and full forms of that
script, and CI invokes the same component commands.

Deterministic tests fingerprint all five public artifacts for MAYO1 and
MAYO2. Negative tests cover statement binding, cross-key use, limit changes,
nonce equality at the bound, rollback duplicates, canonical decoding, and
sampled wire mutations. Fuzz targets cover all seven decoders plus valid-key
authenticated paths. These checks are designed to catch implementation drift;
they do not substitute for the proof and review work in
Section~\ref{sec:security}.

\bibliographystyle{plainnat}
\bibliography{paper/references}

\end{document}
